\subsection{Space and Time}

\begin{tcolorbox}[title=Problem 8, breakable]
    \begin{enumerate}
        \item Use the definition (1.7) to prove that the scalar product 
        is distributive, that is $\vec{r} \cdot (\vec{u} + \vec{v}) = \vec{r} \cdot \vec{u} + \vec{r} \cdot \vec{v}$. 
        \item If $\vec{r}$ and $\vec{s}$ are vectors that depend on time, prove that the product rule for differentiating 
        products applies to $\vec{r} \cdot \vec{s}$, that is, that 
        \[\frac{d}{ds} (\vec{r} \cdot \vec{s}) = \vec{r} \cdot \frac{d \vec{s}}{dt} + \frac{d \vec{r}}{dt} \cdot \vec{s}.\]
    \end{enumerate}
\end{tcolorbox}

\begin{proof}
    We have 
    \[
        \vec{r} \cdot (\vec{u} + \vec{v}) 
        = \sum_{i = 1}^3 r_i (u_i + v_i) 
        = \sum_{i = 1}^3 (r_i u_i + r_i v_i) 
        = \sum_{i = 1}^3 r_i u_i + \sum_{i = 1}^3 r_i v_i 
        = \vec{r} \cdot \vec{u} + \vec{r} \cdot \vec{v}.
    \]
\end{proof}

\begin{proof}
    Suppose $\vec{r}$ and $\vec{s}$ depend on time $t$. Then
    \[
        \frac{d}{dt} (\vec{r} \cdot \vec{s}) 
        = \frac{d}{dt} \left( \sum_{i=1}^3 r_i s_i \right)
        = \sum_{i=1}^3 \frac{d}{dt}(r_i s_i)
        = \sum_{i=1}^3 \left( \frac{d r_i}{dt} s_i + r_i \frac{d s_i}{dt} \right)
        = \frac{d \vec{r}}{dt} \cdot \vec{s} + \vec{r} \cdot \frac{d \vec{s}}{dt}.
    \]
\end{proof}

\newpage
\begin{tcolorbox}[title=Problem 10, breakable]
    A particle moves in a circle (center $O$ and radius $R$)
    with constant angular velocity $\omega$ counterclockwise.
    The circle lies in the $xy$ plane and the particle is on the $x$
    axis at time $t = 0$. Show that the particle's position is given by 
    \[\vec{r}(t) = \hat{x} R \cos(\omega t) + \hat{y} R \sin(\omega t).\]
    Find the particle's velocity and acceleration. 
    What are the magnitude and direction of the acceleration?
    Relate your results to well-known properties of uniform circular motion.
\end{tcolorbox}

\begin{proof}
    We begin with the position of a point $P$ lying on the unit circle 
    \[
        P = (\cos x, \sin x).
    \]
    Now we scale the radius of the circle by $R$ such that the point $P$ lies $R$ away from $O$
    \[
        P = (R\cos x, R\sin x).
    \]
    We let $x = \omega t$ be the position at time $t$ for angular velocity $\omega$, giving 
    \[
        \vec{r}(t) = (R\cos(\omega t), R\sin(\omega t)) 
        = \hat{x} R \cos(\omega t) + \hat{y} R \sin(\omega t).
    \]
    The velocity vector is obtained by differentiating
    \[
        \vec{v}(t) = \frac{d\vec{r}}{dt}
        = (-R\omega \sin(\omega t), R\omega \cos(\omega t)).
    \]
    The acceleration vector is obtained by differentiating again
    \[
        \vec{a}(t) = \frac{d\vec{v}}{dt}
        = (-R\omega^2 \cos(\omega t), -R\omega^2 \sin(\omega t))
        = -\omega^2 \vec{r}(t).
    \]
    The force vector is pulling the point towards the center of the circle.
    The acceleration vector is in the same direction as the force vector,
    so it points toward the origin.
    Thus the magnitude is 
    \[
        \|\vec{a}(t)\| = \omega^2 R,
    \]
    and the direction is toward the center (opposite to the position vector $\vec{r}(t)$).
\end{proof}

\textbf{Solution:}
Imagine you have a string with a ball tied to the end.
As you increase the velocity of the ball you feel a greater force as the ball tries to escape its path.
This corresponds to the $\omega^2$ term in the force magnitude formula.
Now imagine you half the length of the string and spin with equivalent velocity.
You will feel less force corresponding to $\frac{1}{R}$ in the magnitude $\omega^2 R$.

\begin{tcolorbox}[title=Problem 11, breakable]
\end{tcolorbox}

\begin{tcolorbox}[title=Problem 12, breakable]
\end{tcolorbox}

\begin{tcolorbox}[title=Problem 17, breakable]
\end{tcolorbox}